\scanpage{92}
\mspage{10}
\inserted{(où $t:S\longrightarrow U_{03}$)} par la notation \unsure{anodine|ordinaire}
\[
(g,\lambda).\zeta\;,
\]
et on trouve que l'on a, pour l'action modulaire $T_{g,\lambda}$ sur $\zeta$ :
\[
T_{g,\lambda}\zeta=\lambda\deleted{\operatorname{sg}(g)}(g,\lambda).\zeta
\qquad \unsure{\zeta}\in\Gamma\underline O_{\unsure{U_{0,3}|\mathbb P_{0,3}|U_{03}}}(1)(t) \tag*{(51)}
\]
donc
\[
T_{g,\lambda}\deleted{\unsure{\zeta^2|\zeta}}y=\lambda^2(g,\lambda).\deleted{\unsure{\zeta^2}}y
\qquad y\in\Gamma\underline O_{U_{0,3}}(2)(t) \tag*{(52)}
\]
donc
\[
\begin{aligned}
T_{g,\lambda}(\tau)&=\lambda^2\operatorname{sg}(g)\,(g,\lambda).\tau
\qquad \tau\in\Gamma\deleted{\unsure{\underline O_{U_{03}}(t)}}\\
&\hphantom{{}=\lambda^2\operatorname{sg}(g)\,(g,\lambda).\tau\qquad \tau\in\Gamma}
\underline O_{U_{0,3}}(2)\unsure{(\tfrac12)_0|(\tfrac12)}(t)
\end{aligned} \tag*{(53)}
\]

Quand on se restreint au sous-groupe $\mathbb D'_{3\,S_0}$, pour lequel on a $\lambda^2=\operatorname{sg}(g)$, on trouve donc pour (51) (52) (53)
\[
T_{g\deleted{\unsure{,\lambda}}}\deleted{\lambda}\deleted{\unsure{y|\zeta}}\zeta=
\deleted{\operatorname{sg}(g)}\chi_4(g)\,(g.\zeta) \tag*{(51 bis)}
\]
\[
T_g\,y=\operatorname{sg}(g^0)\,(g.y) \tag*{(52 bis)}
\]
\[
T_g(\tau)=\deleted{\unsure{\lambda|1}}g.\tau \tag*{(53 bis)}
\]
où $g\in\Gamma(\mathbb D'_{3\,S})$ pour $S$ sur $S_0$
\[
\chi_4:\mathbb D'_{3\,S_0}\longrightarrow\mu_4
\quad\text{le caractère canonique} \tag*{(54)}
\]
(de sorte que $\mathbb D'_{3\,\mathrm{ab}}\simeq\mu_4$)
% archive sheet number: 91
